% !TEX program = xelatex
% FedCSL paper in NeurIPS format. Requires neurips_2023.sty from:
% https://neurips.cc/Conferences/2023/AuthorInstructions
% 本文含中文，必须使用 XeLaTeX（Overleaf 中请在 Menu -> Compiler 选择 XeLaTeX）。
\documentclass{article}

% 左侧行号：默认投稿模式由 neurips_2023 自动加 lineno。若需关闭行号，改为 \usepackage[preprint]{neurips_2023}；终稿用 [final]
% [nonatbib]：不让模板自带 author-year 版 natbib，改由下面显式载入 numeric 版
\usepackage[nonatbib]{neurips_2023}
% NeurIPS 发表稿（如 FedPCL NeurIPS 2022）使用的数字方括号引用：[1], [2, 3, 4], [5--7]
\usepackage[numbers,sort&compress]{natbib}

% 使用 XeLaTeX/tectonic 编译，去掉 pdflatex 专用的 inputenc / fontenc / CJKutf8
\usepackage{hyperref}
\usepackage{url}
\usepackage{booktabs}
\usepackage{amsfonts}
\usepackage{amsmath}
\usepackage{amssymb}
\usepackage{amsthm}
\usepackage{bm}
\usepackage{nicefrac}
\usepackage{microtype}
\usepackage{xcolor}
\usepackage{graphicx}
\usepackage{multirow}
\usepackage{array}
\usepackage{algorithm}
\usepackage{algpseudocode}
\usepackage{enumitem}
\usepackage{float}

% 中文支持：XeLaTeX + xeCJK + fandol（tectonic 可自动拉取 fandol-fonts）
\usepackage{fontspec}
\usepackage{xeCJK}
% 用 .otf 文件名而不是字体族名，让 tectonic 能从 bundle 按文件解析
\setCJKmainfont[
  BoldFont      = FandolSong-Bold.otf,
  ItalicFont    = FandolKai-Regular.otf,
  BoldItalicFont= FandolKai-Regular.otf
]{FandolSong-Regular.otf}
\setCJKsansfont[BoldFont=FandolHei-Bold.otf]{FandolHei-Regular.otf}
\setCJKmonofont{FandolFang-Regular.otf}

\hypersetup{
    colorlinks=true,
    linkcolor=black,
    citecolor=black,
    urlcolor=blue,
}

\newtheorem{definition}{定义}
\newtheorem{proposition}{命题}
\newtheorem{theorem}{定理}
\newtheorem{remark}{说明}

\title{FedCSL-Spilter: 面向非 IID 多变量时间序列的自适应尺度拼接联邦 Shapelet 表征学习}

\author{%
  作者姓名 \\
  单位名称 \\
  \texttt{email@example.com}
}

\begin{document}

\maketitle

\renewcommand{\abstractname}{摘要}

\begin{abstract}
多变量时间序列广泛存在于人体活动识别、工业监测、医疗诊断与设备运维等场景中。这类数据通常天然分散在不同终端、机构或边缘节点，既难以集中汇聚，又普遍存在标签稀缺、分布异构和隐私敏感等问题。传统集中式时间序列表征学习方法在拥有完整数据时能够取得较好性能，但在联邦环境下往往面临客户端漂移严重、跨客户端表示不一致以及训练效率受限等挑战。现有联邦时间序列方法多侧重监督分类、参数平均或系统调度，而对可解释的多尺度 shapelet 子模型构造、聚合与效率分析仍然不足。

针对上述问题，本文提出一种面向非独立同分布多变量时间序列的联邦 shapelet 学习框架 FedCSL-Spilter。该方法以 TimeCSL/CSL 的多尺度 shapelet 表征学习思想为基础，在联邦训练中引入三类关键机制。首先，在表示学习层面，本文在本地实例级对比学习目标之上，引入局部--全局联合对比学习与联合知识蒸馏机制，使客户端模型在保留本地适应性的同时，与全局语义空间保持一致。其次，在尺度组合层面，本文设计 Spilter 自适应尺度拼接机制：每个客户端基于本地数据的 STFT/ACF 周期感知评分，从 $R$ 个候选尺度中\textbf{按分数降序选取 top-$m$} 个尺度，将这 $m$ 个尺度的 shapelet 特征在表示维度上拼接组合为一个本地子模型进行端到端训练；客户端仅上传被训练尺度的参数，服务器按尺度粒度完成聚合。再次，在多尺度建模层面，本文结合 STFT 与 ACF 构造周期感知的尺度权重，并将该统计信号同时用于损失加权和本地 top-$m$ 尺度选择。

本文进一步从优化与系统角度分析了所提方法的合理性。我们指出，多尺度 shapelet 的参数量与计算量随着尺度长度增长而显著变化，因此不加区分地激活所有尺度或随机选择尺度都会造成训练预算浪费；基于本地数据特征自适应拼接尺度子模型能够在保持表示能力的同时降低通信和计算开销。实验在 HAR、FaceDetection、LSST、Epilepsy 与 SleepEDF 五个多变量时间序列分类数据集上进行，并从总体性能、消融实验、非 IID 强度、收敛趋势和可解释性分析等多个方面验证方法有效性。结果表明，FedCSL-Spilter 在下游分类性能、表示稳定性和系统效率方面均优于多种联邦表示学习基线与异构模型基线。
\end{abstract}

\noindent\textbf{关键词：} 联邦学习；多变量时间序列；Shapelet；多尺度拼接；自适应尺度选择；对比学习；知识蒸馏

\section{引言}

时间序列分析是人工智能和数据挖掘中的重要研究方向，在智能制造、脑电诊断、人体活动识别、金融预测以及工业告警等场景中均具有广泛应用价值。与静态表格数据或单张图像不同，时间序列不仅包含多维观测值，还包含显著的时间依赖、周期结构和跨尺度模式。因此，如何学习具有判别性、可迁移性和可解释性的时间序列表征，一直是该领域的核心问题之一。

近年来，无监督或自监督时间序列表征学习得到快速发展。相比强依赖标签的监督方法，这类方法能够在标签稀缺甚至无标签的条件下，从原始时间序列中挖掘稳定的语义结构。尤其是基于对比学习的方法，通过构造正负样本关系，使模型在特征空间中捕捉时序信号的内在一致性，从而显著提升下游分类、聚类与异常检测性能\cite{ts2vec,tstcc,tfc}；在此基础上，后续工作进一步从掩码建模\cite{simmtm}、时间--频率一致性\cite{tfc}以及软对比关系\cite{softclt}等不同角度扩展了自监督信号，并在通用时序表示层面取得持续提升\cite{timesurl}。TimeCSL/CSL 工作进一步表明，shapelet 作为时间序列分析中具有天然解释性的子序列模式，不仅可以用于任务特定建模，还可以与对比学习结合，形成通用的无监督 shapelet 表征框架\cite{timescsl}。这为后续研究提供了一个非常有价值的方向，即在保持模型可解释性的前提下提升表示学习能力\cite{interpgn,unishape}。

然而，在现实应用中，大量时间序列数据并不适合集中式处理。一方面，医疗、生理监测、设备传感和边缘终端产生的数据通常分散存储于不同客户端；另一方面，这些数据往往具有较强隐私属性，直接上传原始数据会带来隐私泄露和合规风险。联邦学习为此提供了一种自然的解决方案，即只在本地保留原始数据，通过参数或表示级别的信息交换实现多客户端协同训练\cite{fedavg}。为缓解联邦非 IID 场景下的客户端漂移，现有工作普遍在本地目标中加入约束项（如 FedProx 的近端项\cite{fedprox}）或引入模型级对比监督（如 MOON\cite{moon}）。尽管联邦学习在图像、语言和推荐系统中已被广泛研究，但针对多变量时间序列的联邦表示学习，尤其是可解释的 shapelet 表征学习，仍处于相对早期阶段\cite{fedst,fedtsc}。

把集中式的 TimeCSL/CSL 直接迁移到联邦环境并不是一件简单的事情。首先，联邦场景下不同客户端数据往往呈现严重的非独立同分布特性，不同用户、设备和场景对应的数据统计特征、类别分布和时间模式存在明显差异。若仍采用纯粹的局部训练加参数平均，客户端模型很容易过拟合其本地模式，导致全局模型在表示空间中出现漂移和不稳定现象。其次，时间序列本身具有复杂的多尺度结构，不同数据集甚至同一数据集中的不同样本，其主导模式可能出现在不同的时间尺度上。如果在联邦环境下仍对各尺度一视同仁，则会降低 shapelet 学习的针对性。再次，联邦训练中的“模型异构”同样会体现在尺度异构上：某些尺度的参数量更大、滑窗计算更重，另一些尺度则更轻更快。如果所有尺度都被随机选择或平均对待，系统很容易把预算浪费在低效尺度上。

已有研究从不同角度探讨了异构联邦模型如何适应系统资源。HeteroFL 通过宽度裁剪将同一全局模型转化为不同宽度的本地子模型\cite{heterofl}；FjORD 利用 ordered dropout 提取嵌套子模型，并结合自蒸馏提升训练稳定性\cite{fjord}；DepthFL 通过深度裁剪缓解宽度裁剪在聚合时的参数不匹配问题\cite{depthfl}；ScaleFL 同时在宽度与深度两个方向上做资源自适应裁剪\cite{scalefl}；InclusiveFL 和 AdaptiveFL 强调在强资源异构环境下不能放弃弱设备，而要通过大小模型协同和知识蒸馏共享信息\cite{inclusivefl,adaptivefl}。FlexFL 则把 APoZ-guided flexible pruning、adaptive local pruning 与 self-KD 结合起来，在不确定场景下构造更合适的异构子模型\cite{flexfl}。这些工作表明，模型构造和结构裁剪本身就是联邦效率的关键变量，但它们大多围绕通用 CNN/Transformer 展开，尚未显式考虑多尺度 shapelet 的尺度成本差异与周期效用差异。

基于此，本文提出一种面向非 IID 多变量时间序列的联邦 shapelet 学习与 Spilter 自适应尺度拼接框架 FedCSL-Spilter。本文的核心思想是：以 TimeCSL/CSL 的多尺度 shapelet 对比学习框架为骨干，在联邦环境下同时处理”如何学习可解释表示”和”如何按本地数据特征自适应拼接尺度子模型”两个关键问题。具体而言，本文首先保留本地实例对比学习，并引入局部--全局联合对比学习与联合知识蒸馏，使客户端模型在保留本地适应性的同时，与全局语义空间保持一致。随后，本文围绕 Spilter 设计\textbf{本地分布感知 top-$m$} 尺度拼接机制：每个客户端在训练开始前，基于本地数据的 STFT/ACF 周期评分从 $R$ 个候选尺度中选取分数最高的 $m$ 个尺度，将其 shapelet 特征在表示维度上拼接（stitch）为一个本地子模型进行端到端训练。由于各客户端只训练和上传被选中尺度的参数，模型更新不再是全量传输，而是只围绕各节点自主激活的尺度流动，服务器按尺度粒度聚合全局参数。为了让模型更好地适应不同数据集的频谱与周期结构，本文还基于 STFT 和 ACF 设计周期感知加权机制，对不同 shapelet 尺度进行动态调节，并把该统计信号同时用于损失加权和本地 top-$m$ 选择。

综上，本文试图回答以下几个问题：第一，如何将可解释的 shapelet 表征学习扩展到联邦非 IID 环境；第二，如何让各客户端根据本地数据特征自主拼接尺度子模型，而非依赖服务器统一分发；第三，如何通过局部--全局多尺度知识对齐提升全局表示空间的一致性；第四，如何从理论与实验两个层面解释周期感知加权、自适应尺度拼接与分片聚合的有效性。

本文主要贡献概括如下：
\begin{enumerate}[leftmargin=2.5em]
    \item 提出一种面向联邦多变量时间序列的 shapelet 表征学习与 Spilter 自适应尺度拼接一体化框架，在保留 TimeCSL/CSL 可解释性的基础上，实现联邦环境下的协同表征学习。
    \item 设计局部--全局 joint CL/KD 与 multi-scale CL/KD 机制，通过特征层面的软对齐抑制非 IID 场景下的客户端漂移，并增强不同尺度子空间的一致性。
    \item 引入基于 STFT 与 ACF 的周期感知多尺度权重机制，使 shapelet 学习能够根据不同数据的主导周期结构进行自适应调整；并将该周期信号作为客户端自主选择尺度、拼接子模型的决策依据。
    \item 在尺度组合层面，提出\textbf{纯本地分布感知 top-$m$} 的 Spilter 尺度拼接机制：各客户端仅依据 STFT/ACF 周期评分选取 top-$m$ 尺度并拼接子模型，与随机 top-$m$、固定尺度及全尺度 FedCSL 等策略进行比较。
    \item 从优化与系统效率视角给出方法合理性的理论解释，并在 HAR、FaceDetection、LSST、Epilepsy 与 SleepEDF 数据集上通过总体性能、消融实验、收敛分析和可解释性实验验证所提方法的有效性。
\end{enumerate}

\section{相关工作}

\subsection{时间序列无监督表示学习}

时间序列无监督表示学习旨在在缺乏标注数据的前提下学习通用时序表示。早期方法通常依赖重构任务、序列预测或自编码器结构，但这类方法往往更关注局部重建误差，而不一定能够直接形成适合下游任务的判别特征\cite{simmtm}。随后，对比学习被广泛引入时间序列分析领域，其基本思想是在不同增强视图之间最大化正样本一致性，同时抑制负样本相似性，从而学习稳定而有区分度的表示。典型工作如 TS2Vec 在分层上下文视图上进行对比\cite{ts2vec}，TS-TCC 分别构造时间对比与上下文对比模块\cite{tstcc}，TF-C 将时间与频率视图的一致性作为自监督信号\cite{tfc}；更进一步，SoftCLT 以软分配缓解了“邻近样本被强制推开”的问题\cite{softclt}，TimesURL 则系统性地重审了正负对构造与 SSCL 损失设计\cite{timesurl}，而 SimMTM 通过掩码时序建模与流形邻居重构扩展了自监督范式\cite{simmtm}。

现有大量方法以 CNN、Transformer 或其变体为主干进行时序建模\cite{fcnbaseline,timesnet}，但这些架构最初并非专为时间序列设计，在处理不同长度、多尺度以及强周期特征时常常存在解释性不足的问题。一条互补的路线是以 shapelet 等可解释子序列模式为核心构建时序分类器\cite{interpgn,unishape}。TimeCSL/CSL 的意义在于将 shapelet 这一对时间序列更自然、更可解释的模式表示与对比学习相结合，使模型能够从多个尺度和多种相似性度量出发提取通用特征\cite{timescsl}。相比纯黑盒表示，shapelet-based representation 具有更强的可解释性，更便于对模型输出进行可视化分析和任务迁移\cite{interpgn}。

\subsection{联邦时间序列学习与知识蒸馏}

联邦时间序列学习研究的核心问题在于：在不共享原始数据的条件下，如何使多个客户端协同训练得到兼顾隐私与性能的模型\cite{fedavg}。现有工作多聚焦于监督分类场景，其中参数平均、个性化联邦学习、蒸馏和迁移学习是几种常见技术路线\cite{fedkdsurvey}。在时序领域，FedST 提出面向 shapelet 变换的安全联邦协议\cite{fedst}，FedTSC 进一步将可解释 TSC 模型嵌入联邦系统\cite{fedtsc}；EFDLS 则将联邦蒸馏引入多任务时间序列分类，通过特征级教师--学生蒸馏和权重匹配策略促进不同任务间知识共享\cite{efdls}。这类工作说明，在时序数据异构较强的环境中，仅靠参数平均并不足以支持有效的跨客户端知识传递，特征层面的蒸馏和语义对齐具有重要意义\cite{moon,fedx}。

跨领域的联邦无监督表示学习也在快速发展。FedCA 最早针对 FURL 提出跨客户端表示对齐以缓解表示不一致\cite{fedca}；FedMoCo 把 MoCo 迁移到医疗图像联邦场景，辅以元数据迁移与自适应聚合\cite{fedmoco}；FedX 通过双向知识蒸馏与局部--全局对比学习学习无偏表示\cite{fedx}；FedU2 指出 FUSL 同时存在表示坍缩和表示空间不一致，给出统一聚合方案\cite{fedu2}。在有监督侧，MOON 引入模型级对比以纠正本地训练\cite{moon}，FedProc 与 FedPCL 以类原型构造对比信号\cite{fedproc,fedpcl}，Reducing non-IID (CDL) 则用对比散度损失抑制跨车辆异构\cite{fedcdl}，并有工作从互信息视角统一讨论联邦对比的合理性\cite{fedmi}。更贴近本文局部--全局蒸馏思路的是 FLESD\cite{flesd}：它利用集成相似度矩阵在服务器侧进行相似度蒸馏，从而在强约束下实现高效联邦表示学习；近期工作 FLPD\cite{flpd} 与 FedRCIL\cite{fedrcil} 则分别从原型相似度蒸馏与增量对比学习出发，给出通信友好的联邦蒸馏方案。此外，FedCOG\cite{fedcog}、MrTF\cite{mrtf}、FedIOCD\cite{fediocd} 与 FedNoRo\cite{fednoro} 等工作从“生成对齐”“转导推理”“输入输出协同蒸馏”和“噪声鲁棒聚合”等角度进一步拓展了联邦蒸馏的边界。

然而，上述方法大多默认任务目标为监督分类或视觉无监督场景，很少直接研究无监督 shapelet 表征学习问题。相比之下，本文更加关注联邦环境下的通用时序表示，即先学习可解释的表示空间，再在该空间上完成下游任务建模。这一思路既继承了 TimeCSL 的优势\cite{timescsl}，也更适合标签稀缺场景\cite{fedsmall}。

\subsection{异构联邦模型与尺度级自适应组合}

在真实联邦环境中，由于带宽、设备异构性和大规模客户端数量的限制，服务器通常无法假设所有客户端都能承担同样规模的模型训练。Oort 和 FedCS 类方法表明，统计效用、设备速度、带宽与训练时延共同决定了联邦系统的 time-to-accuracy 表现\cite{oort,fedcsadaptive}。本文不再展开客户端采样问题，而是沿用这种系统效用视角，进一步关注在给定客户端参与集合下，如何让各节点根据本地数据特征自主组合出适合自身的尺度子模型。

异构模型联邦通常从同一个全局模型中派生多个子模型。HeteroFL 通过按宽度裁剪通道数得到不同复杂度的子模型\cite{heterofl}；DepthFL 沿深度方向构造不同层数的子网络\cite{depthfl}；FjORD 使用 ordered dropout 得到嵌套子模型，从而让不同容量的客户端共享一部分有序参数\cite{fjord}；ScaleFL 同时考虑宽度与深度比例，使子模型能更细致地匹配客户端资源\cite{scalefl}；InclusiveFL 与 AdaptiveFL 则从弱设备参与和资源约束角度强调大模型、小模型之间的协同训练与蒸馏\cite{inclusivefl,adaptivefl}。FlexFL 进一步利用 APoZ 估计层内神经元冗余，并让设备根据当前资源自适应裁剪收到的模型，同时通过 self-KD 改善大模型性能\cite{flexfl}。

上述工作说明，联邦系统效率不只取决于客户端是否参与，也取决于子模型的结构粒度能否与设备资源相匹配。与 CNN 中”通道/层”的裁剪不同，FedCSL 的天然结构单元是 shapelet 尺度：每个尺度对应一组 shapelets 及其匹配算子。更重要的是，不同客户端因本地数据的周期/频谱结构不同，对不同尺度的需求天然存在差异。因此，本文将异构模型构造具体化为\textbf{客户端自主尺度拼接}：各客户端不再被动接收服务器指派的统一全尺度模型，而是\textbf{仅依据本地周期感知评分}从 $R$ 个候选尺度中选取 top-$m$ 个尺度，将其 shapelet 特征在表示维度上拼接（stitch）为一个本地子模型进行训练，仅上传被训练尺度的参数；设备显存/算力差异则通过为不同客户端配置不同的 $m$ 来体现（见 §\ref{sec:resource-hetero}），而不参与尺度排序本身。

\section{问题定义与总体框架}

\subsection{联邦时序场景定义}

设联邦系统由一个中心服务器和 $K$ 个客户端构成。第 $k$ 个客户端持有本地多变量时间序列数据集
\begin{equation}
D_k = \{x_{k,i}\}_{i=1}^{n_k},
\end{equation}
其中 $x_{k,i}\in \mathbb{R}^{C\times T}$，$C$ 表示通道数，$T$ 表示时间长度，$n_k$ 表示客户端 $k$ 的样本规模。总样本规模为
\begin{equation}
N = \sum_{k=1}^{K} n_k.
\end{equation}

在本文中，不同客户端之间的数据分布通常满足非独立同分布条件，即存在
\begin{equation}
P_k(x)\neq P_j(x),\quad k\neq j.
\end{equation}
这种差异不仅体现在类别分布上，也体现在时间尺度、周期模式、幅值范围以及局部噪声结构上。

\subsection{目标定义}

本文的目标不是直接训练一个端到端分类器，而是学习一个联邦共享的 shapelet 编码器 $f(\cdot;w)$，将原始时间序列映射到可解释的表示空间：
\begin{equation}
z = f(x;w)\in \mathbb{R}^{d}.
\end{equation}
得到表示后，再通过下游分类器（如 SVM）完成分类任务评估。该设计与 TimeCSL 的统一分析范式保持一致，即将表示学习与任务求解解耦。

\subsection{总体框架}

FedCSL-Spilter 的总体流程如图~\ref{fig:framework} 所示，可概括为以下步骤：
\begin{enumerate}[leftmargin=2.5em]
    \item 服务器初始化全局 shapelet 编码器，各客户端基于本地数据预计算周期感知尺度效用分数；
    \item 每个客户端根据本地周期感知评分按 top-$m$ 选取 $m$ 个尺度，将这 $m$ 个尺度的特征拼接为本地子模型；
    \item 客户端从服务器拉取被选中尺度的全局参数，在本地对拼接后的子模型进行 shapelet 对比学习，并引入局部--全局联合蒸馏与尺度级蒸馏；
    \item 各客户端仅上传被训练尺度的参数更新；
    \item 服务器按尺度收集各客户端上传的参数，对每个尺度分别执行样本量加权聚合，未被本轮上传的尺度保持服务器原值；
    \item 周期感知模块根据 STFT 与 ACF 统计动态调整多尺度损失权重，并为下一轮本地尺度选择提供效用信号；
    \item 重复上述过程直到训练收敛，并在学习到的全局表示上执行下游评估。
\end{enumerate}

\begin{figure}[H] % 放在正文引用位置附近
    \centering
    \includegraphics[width=0.98\linewidth]{imageall.png}
    \caption{FedCSL-Spilter 整体框架图。上半部分为服务器端：Scale-wise Aggregator 按尺度聚合各客户端上传的参数更新，得到全局 Shapelet 编码器 $w_g$，下游以 SVM 做表示评测。下半部分为客户端 $k$ 的本地训练：多变量时间序列经数据增强形成两种视图 $x_q,x_k$；客户端根据本地周期特征自主选出 $m$ 个尺度，将其 shapelet 特征拼接（stitch）为本地子模型；特征对齐节点聚合本地与全局特征，驱动局部对比损失 $\mathcal{L}_{local}$、局部--全局联合对比 $\mathcal{L}_{joint\text{-}CL}$ 与联合蒸馏 $\mathcal{L}_{joint\text{-}KD}$；右下 Period-Aware Local Weighting Module 由 STFT（短--中尺度）与 ACF（中--长尺度）两分支产生尺度权重 $\bm{\pi}$，同时反馈到多尺度损失和本地尺度选择。}
    \label{fig:framework}
\end{figure}


\section{方法}

\subsection{多尺度 Shapelet 编码器}

本文采用与 TimeCSL/CSL 一致的多尺度 shapelet 表示思想。对于长度为 $T$ 的输入时间序列，在 $\left[0.1T,0.8T\right]$ 范围内均匀选取 $R$ 个尺度，每个尺度学习固定数量的 shapelets。不同于只依赖单一相似性度量的编码方式，本文实现同时支持欧氏距离、余弦相似度、互相关以及它们的混合形式。记第 $r$ 个尺度下的 shapelet 集合为 $\mathcal{S}^{(r)}=\{s_m^{(r)}\}_{m=1}^{M_r}$，则一个样本的最终表示可写为
\begin{equation}
z = \left[z^{(1)},z^{(2)},\dots,z^{(R)}\right],
\end{equation}
其中 $z^{(r)}$ 表示输入序列与尺度 $r$ 下所有 shapelets 的最优匹配响应向量。

这种表示方式有两个优势。其一，不同尺度的 shapelets 分别捕捉局部瞬时模式、中程结构和长期依赖，使模型能覆盖更广的时序模式空间。其二，shapelet 作为具体的子序列模式，天然具备解释性，便于后续通过可视化分析理解模型决策依据。

\subsection{本地实例级对比学习}

对每个被选中的客户端，本文首先在本地样本上构造两种增强视图。增强操作包括加噪、裁剪、池化、量化和时间扭曲等时序特有变换，这一构造延续了 TS2Vec/TS-TCC 等时序对比方法的做法\cite{ts2vec,tstcc}。给定输入样本 $x$，记两种增强视图为 $x_q$ 和 $x_k$，对应表示分别为
\begin{equation}
q = f(x_q;w_k),\qquad k = f(x_k;w_k).
\end{equation}
对表示进行归一化后，构造实例级对比学习目标：
\begin{equation}
\mathcal{L}_{\text{local}} = \mathrm{CE}\left(\frac{qk^\top}{\tau}, y\right),
\end{equation}
其中 $\tau$ 为温度系数，$y$ 表示正样本匹配标签。该目标鼓励同一样本的不同增强视图在表示空间中接近，同时与其他样本保持区分。

\subsection{局部--全局联合对比与知识蒸馏}

在联邦非 IID 条件下，仅依赖本地实例级对比学习容易导致各客户端在不同语义坐标系中优化，进而削弱全局模型聚合效果\cite{fedu2}。这一问题在视觉联邦无监督工作中已被广泛讨论：MOON 通过模型级对比把本地表示拉回全局邻域\cite{moon}，FedX 引入跨客户端双向蒸馏学习无偏表示\cite{fedx}，FLESD 则在服务器侧利用集成相似度矩阵完成相似度蒸馏\cite{flesd}。受此启发，本文将上一轮全局模型作为语义教师，记其在两种增强视图上的表示为
\begin{equation}
q_g = f(x_q; w_g),\qquad k_g = f(x_k; w_g).
\end{equation}

在此基础上，引入联合对比项
\begin{equation}
\mathcal{L}_{\text{joint-CL}} = \mathrm{CE}\left(\frac{q_gk^\top}{\tau}, y\right),
\end{equation}
以及联合蒸馏项
\begin{equation}
\mathcal{L}_{\text{joint-KD}} = D_{\mathrm{KL}}(q \,\|\, q_g) + D_{\mathrm{KL}}(k \,\|\, k_g),
\end{equation}
其中 $D_{\mathrm{KL}}$ 表示基于相似度矩阵构造的 KL 散度。通过这两项，本地模型既保留了对本地时序模式的适应性，又被限制在全局共享语义空间附近，从而减小由数据异构引起的客户端漂移。

\subsection{多尺度结构约束与尺度级蒸馏}

对于表示向量 $z=[z^{(1)},\dots,z^{(R)}]$，整体级别的蒸馏并不能保证每个尺度子空间都学得充分。为此，本文将联合对齐细化到每个尺度：
\begin{equation}
\mathcal{L}_{\text{scale-CL}} = \sum_{r=1}^{R}\alpha_r \,
\mathrm{CE}\left(\frac{q_g^{(r)}(k^{(r)})^\top}{\tau}, y\right),
\label{eq:scale-cl}
\end{equation}
\begin{equation}
\mathcal{L}_{\text{scale-KD}} = \sum_{r=1}^{R}\alpha_r \left[
D_{\mathrm{KL}}\left(q^{(r)} \,\|\, q_g^{(r)}\right)+
D_{\mathrm{KL}}\left(k^{(r)} \,\|\, k_g^{(r)}\right)\right].
\end{equation}
其中 $\alpha_r$ 为第 $r$ 个 shapelet 尺度上的损失权重。启用周期感知（\texttt{UseACF}）时，$\alpha_r$ 由第~\ref{sec:period-score} 节在线计算；否则可取均匀权重。

此外，本文保留了来自 CSL/TimeCSL 的结构性约束。设某个尺度下的表示矩阵为 $Z^{(r)}$，其协方差矩阵为
\begin{equation}
C^{(r)} = \frac{\left(Z^{(r)}\right)^\top Z^{(r)}}{B-1},
\end{equation}
则通过约束协方差非对角项，可以抑制不同 shapelet 通道之间的冗余，提升各尺度表示的独立性与判别性。结合多尺度累积统计，本文形成了结构对齐损失 $\mathcal{L}_{\text{structure}}$，以避免表示坍缩和尺度冗余。

\subsection{周期感知加权机制}
\label{sec:period-score}

多变量时间序列在不同数据集、不同通道上的\textbf{主导周期/主频}往往不同；而 shapelet 编码器在 $R$ 个离散长度上并行提取模式。若对所有尺度等权优化，容易在\textbf{与当前数据频谱/周期不匹配}的尺度上浪费梯度。本文在实现中于每个训练步、基于\textbf{当前 batch} 的原始输入 $\mathcal{X}\in\mathbb{R}^{N\times K\times L}$（$N$ 为 batch 大小，$K$ 为通道数，$L$ 为序列长度）调用 \texttt{period\_score}（见 \texttt{utils.py}），构造与 $R$ 个参考尺度对齐的非负权重，再写入式~\eqref{eq:scale-cl} 及尺度级蒸馏中的 $\alpha_r$（代码中另乘常数以调节梯度幅度，不影响相对比例）。

\paragraph{参考尺度与离散锚点。}
取归一化锚点
\begin{equation}
\rho_r=\rho_{\min}+\frac{r-1}{R-1}\left(\rho_{\max}-\rho_{\min}\right),
\qquad r=1,\ldots,R,
\end{equation}
其中实现中采用 $\rho_{\min}=0.1,\rho_{\max}=0.8$。
对应以采样点为单位的\textbf{目标尺度}（实现中对每个样本独立计算）
\begin{equation}
\tau_r=\bigl\lfloor \rho_r L \bigr\rfloor,\qquad r=1,\ldots,R.
\end{equation}
这些锚点与多尺度 shapelet 的长度配置一一对应。

\paragraph{基于 STFT 的频域得分（短—中尺度分支）。}
对 batch 内每个样本的至多 $k=\min(9,K)$ 条通道分别计算短时傅里叶变换。实现采用 \texttt{scipy.signal.stft}，采样率归一化为 $f_s{=}1$，窗长 $W=\max(4,\lfloor L/2\rfloor)$，FFT 长度 $N_{\mathrm{fft}}=8W$，帧移为 $h=\lfloor W/4\rfloor$，相邻帧重叠 $W{-}h$。记通道 $c$ 的功率谱为 $P_c(f,t)=|Z_{xx}(f,t)|^2$。为突出\textbf{能量集中频带}，先取功率的 $95\%$ 分位数作为阈值，构造高能量掩膜 $\mathbb{I}[P_c\ge Q_{0.95}(P_c)]$，再定义\textbf{掩膜加权平均频率}
\begin{equation}
\bar{f}_c=\frac{\sum_{f,t} f\cdot P_c(f,t)\cdot \mathbb{I}[\cdot]}{\sum_{f,t} P_c(f,t)\cdot \mathbb{I}[\cdot]+\varepsilon}.
\end{equation}
在 $f_s{=}1$ 的约定下，$\bar{f}_c$ 可视为该通道\textbf{主频}的代理；实现中将其映射为与“周期长度”同量纲的标量 $s_c^{\mathrm{ST}}=1/(\bar{f}_c+\varepsilon)$（$\bar{f}_c$ 过小则置 $0$）。随后用高斯核将标量 $s_c^{\mathrm{ST}}$ 软对齐到各锚点（实现中对 $r=1,\ldots,R$ 均计算，再对通道在 $r$ 上归一化）：
\begin{equation}
u_{c,r}^{\mathrm{ST}}=\exp\!\left(-\frac{\bigl(s_c^{\mathrm{ST}}-\tau_r\bigr)^2}{2\sigma^2}\right),\quad
\sigma=0.1\,L,\quad r=1,\ldots,R,
\end{equation}
对 batch 与通道累加后仅保留短--中尺度索引集合 $\mathcal{I}_{\mathrm{ST}}=\{1,\ldots,\lfloor R/2\rfloor\}$ 上的得分 $\mathbf{S}^{\mathrm{ST}}$，与代码中的短尺度分支一致。图~\ref{fig:period-stft} 预留为\textbf{单通道 STFT 功率谱/时频图}示意，可标注高能量区域与 $\bar{f}_c$ 的几何含义。


\begin{figure}[htbp]
    \centering
    \includegraphics[width=0.9\linewidth]{../../figs/stft_v5_paper_branch.png}
    \caption{周期感知中 STFT 分支的频谱示意（HAR 真实样本：Walking \#2759 / total\_acc\_x）：左图为单通道 STFT 功率谱，白色等高线为 $95\%$ 高能量掩膜，黄色虚线为加权主频 $\bar f_c$；右图将标量 $1/\bar f_c$ 通过高斯核软对齐到 $R$ 个 shapelet 尺度锚点 $\tau_r$，其中 $\mathcal{I}_{\mathrm{ST}}$ 中的短--中尺度由 STFT 分支写入 $\mathbf{S}^{\mathrm{ST}}$，其余中长--长尺度由 ACF 分支接管。}
    \label{fig:period-stft}
\end{figure}

\paragraph{基于 ACF 的时域滞后得分（中长—长尺度分支）。}
对同一批通道，使用 \texttt{statsmodels} 的样本自相关函数（\texttt{fft=True}，最大滞后 $L{-}1$）得到 $R_c(\ell)$，$\ell=0,\ldots,L{-}1$。实现仅在\textbf{后半段滞后} $\ell\ge \lfloor L/2\rfloor$ 上，用 \texttt{scipy.find\_peaks} 检测峰值，且要求 $R_c(\ell)>\theta_{\mathrm{acf}}$（代码中 $\theta_{\mathrm{acf}}=0.2$）；方差过小的常数序列跳过。每个峰值滞后 $\ell$ 对全部 $R$ 个锚点产生反距离投票
\begin{equation}
v_{c,r}(\ell)=\frac{1}{|\tau_r-\ell|+\varepsilon},\qquad
\mathbf{v}_c=\sum_{\ell\in\mathcal{P}_c}\bigl(v_{c,1}(\ell),\ldots,v_{c,R}(\ell)\bigr),
\end{equation}
其中 $\mathcal{P}_c$ 为通道 $c$ 上满足条件的峰值集合；再对 $\mathbf{v}_c$ 按通道归一化并在 batch 与通道上累加，得到 $\mathbf{S}^{\mathrm{ACF}}\in\mathbb{R}^R$。实现中仅取中长--长尺度索引集合 $\mathcal{I}_{\mathrm{ACF}}=\{\lfloor R/2\rfloor+1,\ldots,R\}$ 上的得分，与较长 shapelet 尺度对齐，以强调\textbf{长周期/长记忆}结构。图~\ref{fig:period-acf} 给出\textbf{ACF 曲线随滞后 $\ell$} 的示意，标出阈值 $\theta_{\mathrm{acf}}$ 与后半段峰值。

\begin{figure}[htbp]
    \centering
    \includegraphics[width=0.9\linewidth]{../../figs/acf_v5_paper_branch.png}
    \caption{周期感知中 ACF 分支的滞后—投票示意：左图为样本自相关函数 $R_c(\ell)$ 随滞后 $\ell$ 的变化，红色虚线为阈值 $\theta_{\mathrm{acf}}$，后半段（$\ell\ge\lfloor L/2\rfloor$）检出的峰值以圆点标记；右图将这些峰值按反距离投票软对齐到 $R$ 个 shapelet 尺度锚点 $\tau_r$，$\mathcal{I}_{\mathrm{ACF}}$ 中的中长--长尺度由 ACF 分支写入，并与 STFT 分支互补。}
    \label{fig:period-acf}
\end{figure}

\paragraph{双分支融合与损失中的使用。}
将 STFT 分支向量归一化得 $\hat{\mathbf{S}}^{\mathrm{ST}}=\mathbf{S}^{\mathrm{ST}}/\|\mathbf{S}^{\mathrm{ST}}\|_1$，ACF 分支在 $\mathcal{I}_{\mathrm{ACF}}$ 上的向量归一化得 $\hat{\mathbf{S}}^{\mathrm{ACF}}_{\mathcal{I}_{\mathrm{ACF}}}$。融合系数 $\beta\in[0,1]$（实现中对应配置项 \texttt{beta}，调用 \texttt{period\_score(..., alpha=self.beta)}）控制频域与滞后域信任比例：
\begin{equation}
\bm{\pi}=\mathrm{Norm}\!\left(
\beta\cdot \mathrm{Embed}_{\mathcal{I}_{\mathrm{ST}}}\!\bigl(\hat{\mathbf{S}}^{\mathrm{ST}}\bigr)
\;+\;
(1-\beta)\cdot \mathrm{Embed}_{\mathcal{I}_{\mathrm{ACF}}}\!\bigl(\hat{\mathbf{S}}^{\mathrm{ACF}}_{\mathcal{I}_{\mathrm{ACF}}}\bigr)
\right),
\label{eq:period-fuse}
\end{equation}
其中 $\mathrm{Embed}_{\mathcal{I}}(\cdot)$ 表示把子向量嵌回 $R$ 维尺度空间、未覆盖尺度补零，$\mathrm{Norm}(\cdot)$ 表示对 $R$ 维向量做非负归一化使和为 $1$；若各分支全零则退化为均匀权重。实现返回前将 $\bm{\pi}$ 乘以常数 $10$ 以便与损失量级匹配；在 \texttt{train.py} 中，第 $r$ 个尺度的对比/蒸馏项再乘以与 $\pi_r$ 成正比的系数（代码为 $\pi_r\times 5$ 写入 \texttt{precisions}）。因此，\textbf{同一前向 batch} 的频谱与 ACF 统计直接决定该步多尺度损失的相对 emphasis，无需额外标签。

\subsection{Spilter 自适应尺度拼接与本地分布感知 Top-$m$ 选择}
\label{sec:spilter}

多尺度 shapelet 编码器天然可以拆分为 $R$ 个尺度子模型。记第 $r$ 个尺度的参数为 $\theta^{(r)}$，尺度长度为 $\ell_r$，shapelet 数为 $M_r$。完整模型参数可写为
\begin{equation}
w=\{\theta^{(1)},\theta^{(2)},\ldots,\theta^{(R)},\theta^{(\mathrm{shared})}\},
\end{equation}
其中 $\theta^{(\mathrm{shared})}$ 表示投影、归一化或分类头等非尺度参数。与现有异构联邦工作中服务器统一规划子模型结构不同，Spilter 的核心理念是：\textbf{每个客户端根据本地数据的周期/频谱分布感知评分，自主决定哪些尺度对自身最有价值，并将这些尺度的 shapelet 特征拼接组合为一个本地子模型进行训练}。具体而言，客户端 $k$ 从 $R$ 个候选尺度中自主选出 $m$ 个尺度（$m<R$），构成激活尺度集合 $\mathcal{R}_k\subseteq\{1,\ldots,R\}$，仅从服务器拉取并训练 $\{\theta^{(r)}:r\in\mathcal{R}_k\}$ 的参数。

\paragraph{本地周期感知尺度评分。}
客户端 $k$ 首先基于本地数据，利用第~\ref{sec:period-score} 节的 STFT/ACF 双分支机制计算周期感知尺度分数向量 $\mathbf{s}_k=(s_{k,1},\ldots,s_{k,R})$。该分数反映了客户端的本地周期/频谱结构与各 shapelet 尺度之间的匹配程度：若客户端数据的主导周期与某尺度的锚点 $\tau_r$ 接近，则该尺度获得较高分数。与服务器集中式计算不同，该评分过程完全在客户端本地完成，无需上传原始数据，既保护了数据隐私，又使尺度选择决策更贴合本地数据分布。

\paragraph{本地分布感知 Top-$m$ 选择。}
获得 $\mathbf{s}_k$ 后，客户端 $k$ 直接按分数降序选取前 $m$ 个尺度构成激活集合，不再引入参数量、滑窗计算量或通信成本等系统侧代理项：
\begin{equation}
\mathcal{R}_k=\mathrm{TopM}(\mathbf{s}_k,m)
=\underset{\substack{\mathcal{R}\subseteq\{1,\ldots,R\}\\|\mathcal{R}|=m}}{\arg\max}\;
\sum_{r\in\mathcal{R}} s_{k,r}.
\label{eq:spilter-topm}
\end{equation}
实现中对 $\mathbf{s}_k$ 做非负归一化后以稳定排序选取 top-$m$；同分情况下优先选择较小尺度索引（\texttt{local\_score\_topm} 模式）。该计划在训练开始前基于客户端本地数据统计一次性生成并保持固定。当 $m=R$ 时退化为全尺度 FedCSL；若将 $s_{k,r}$ 替换为均匀随机数再取 top-$m$，则得到 Random-top-$m$ 消融版本。需要强调的是，\textbf{尺度排序完全由本地分布感知评分驱动}；不同客户端因数据周期结构不同而自然得到不同的 $\mathcal{R}_k$，从而在全局层面形成多视图集成，而无需在评分函数中显式写入系统成本项。

\paragraph{尺度特征拼接与子模型训练。}
选定激活尺度集合 $\mathcal{R}_k$ 后，客户端从服务器拉取对应尺度的全局参数，将 $m$ 个尺度的 shapelet 特征在表示维度上进行拼接（stitch），构造本地子模型。具体而言，对于输入样本 $x$，客户端通过编码器分别提取每个激活尺度的特征 $z^{(r)} = f^{(r)}(x; \theta^{(r)})$，$r\in\mathcal{R}_k$，再将它们沿特征维度拼接为统一的子模型表示：
\begin{equation}
z^{\mathrm{stitch}}_k = \bigl[z^{(r_1)} \,\|\, z^{(r_2)} \,\|\, \cdots \,\|\, z^{(r_m)}\bigr],
\qquad r_i\in\mathcal{R}_k .
\label{eq:stitch}
\end{equation}
该拼接表示随后被用于计算本地实例级对比损失 $\mathcal{L}_{\text{local}}$，以及局部--全局联合对比项 $\mathcal{L}_{\text{joint-CL}}$ 与联合蒸馏项 $\mathcal{L}_{\text{joint-KD}}$。此外，为了保持各尺度子空间的语义一致性，每个激活尺度还独立计算尺度级对比损失 $\mathcal{L}_{\text{scale-CL}}$ 和尺度级蒸馏损失 $\mathcal{L}_{\text{scale-KD}}$，并以周期感知权重 $\alpha_r$ 加权。换言之，客户端训练的是一个由 $m$ 个尺度拼接而成的”定制化子模型”，其主干损失在拼接表示上计算，辅助对齐损失在各分量尺度上分别计算，从而在保持整体表示质量的同时，确保每个分量尺度与全局语义空间对齐。

\subsection{分片通信与尺度级聚合}

Spilter 不上传完整客户端模型，而是按尺度分片通信。第 $t$ 轮，客户端 $k$ 根据本地确定的激活尺度集合 $\mathcal{R}_k$，从服务器拉取对应尺度参数：
\begin{equation}
\Theta_{k}^{\downarrow}(t)=\{\theta_g^{(r,t)}:r\in\mathcal{R}_k\},
\end{equation}
本地训练后仅上传被训练尺度的参数更新：
\begin{equation}
\Theta_{k}^{\uparrow}(t)=\{\theta_k^{(r,t+1)}:r\in\mathcal{R}_k\}.
\end{equation}
因此一轮双向通信量为
\begin{equation}
\mathcal{C}_{\mathrm{split}}(t)
=2\sum_{k=1}^{K}\sum_{r\in\mathcal{R}_k}B_r + O(KR),
\end{equation}
其中 $O(KR)$ 为尺度索引元数据，远小于参数传输。若所有客户端训练完整 $R$ 尺度模型，则通信量为
\begin{equation}
\mathcal{C}_{\mathrm{full}}(t)
=2K\sum_{r=1}^{R}B_r.
\end{equation}
当各尺度参数量近似相等时，$\mathcal{C}_{\mathrm{split}}/\mathcal{C}_{\mathrm{full}}\approx m/R$；不同客户端因本地评分不同而激活不同尺度子集，进一步使全局尺度覆盖更分散。

计算量同理可写为
\begin{equation}
\mathcal{F}_{\mathrm{split}}(t)
=\sum_{k=1}^{K}E\,n_k\sum_{r\in\mathcal{R}_k}F_r,
\qquad
\mathcal{F}_{\mathrm{full}}(t)
=\sum_{k=1}^{K}E\,n_k\sum_{r=1}^{R}F_r,
\end{equation}
其中 $E$ 为本地 epoch 数。由于 $F_r$ 随 $\ell_r$ 和滑窗匹配次数变化，计算节省并不只是尺度数目的一半，而取决于被选择尺度的成本组合。

服务器端按尺度聚合。设
\begin{equation}
\mathcal{K}_r(t)=\{k:r\in\mathcal{R}_k\}
\end{equation}
为本轮上传尺度 $r$ 的客户端集合，则
\begin{equation}
\theta_g^{(r,t+1)}=
\begin{cases}
\displaystyle
\sum_{k\in\mathcal{K}_r(t)}
\frac{n_k}{\sum_{j\in\mathcal{K}_r(t)}n_j}
\theta_k^{(r,t+1)}, & \mathcal{K}_r(t)\neq\emptyset,\\[2ex]
\theta_g^{(r,t)}, & \mathcal{K}_r(t)=\emptyset.
\end{cases}
\label{eq:scale-agg}
\end{equation}
也就是说，某个尺度只有在本轮被至少一个客户端训练并上传时才会更新；否则保持服务器原值。该策略与 FlexFL、HeteroFL 等异构模型联邦的核心思想一致：不同客户端可训练不同规模的子模型，但服务器仍维护一个完整的全局模型\cite{heterofl,flexfl}。

\subsection{总体优化目标}

综上，本文在客户端 $k$ 上的总体目标可写为
\begin{equation}
\begin{aligned}
\mathcal{L}_k ={}& \gamma\,\mathcal{L}_{\text{local}}
+ \lambda_5\,\mathcal{L}_{\text{structure}} \\
&+ \mathbb{I}_{\text{FedCSL}}\Bigl[
\lambda_1\,\mathcal{L}_{\text{joint-CL}}
+ \lambda_2\,\mathcal{L}_{\text{joint-KD}}
+ \lambda_3\,\mathcal{L}_{\text{scale-CL}}
+ \lambda_4\,\mathcal{L}_{\text{scale-KD}}
\Bigr],
\end{aligned}
\end{equation}
其中 $\mathbb{I}_{\text{FedCSL}}$ 在 FedCSL 及其 one-hot/splitteacher 变体中取 $1$，否则取 $0$，其后各项分别由配置项 \texttt{UseJointCL}、\texttt{UseJointKD}、\texttt{UseScaleCL}、\texttt{UseScaleKD} 控制是否计入。$\mathcal{L}_{\text{structure}}$ 在全尺度分支中对应 $\texttt{l3}\times(\mathcal{L}_{\text{cca}}+\texttt{l4}\times\mathcal{L}_{\text{sdl}})$，故式中 $\lambda_5$、$\lambda_{\text{sdl}}$ 与超参 \texttt{l3}、\texttt{l4} 相当；在 Spilter 分支中，为避免额外运行未激活尺度，该结构项可退化为只在被选尺度上计算或暂时关闭。$\lambda_1,\ldots,\lambda_4$ \textbf{不是}互不相关的自由超参：\texttt{train.py} 的 \texttt{update\_CL} 用 $\gamma$、$\zeta=1-\gamma$ 及固定因子 $\frac12$、$\frac{1}{10}$ 对联合项与多尺度项做缩放，且多尺度上 \texttt{UseACF} 时 $\alpha_r$ 与 \texttt{precisions}（含因子 $5$）相乘。各项损失是否开启，由实验配置文件决定。

\begin{algorithm}[H]
\caption{FedCSL-Spilter 训练流程}
\label{alg:fedcsl}
\begin{algorithmic}[1]
\State 初始化全局模型 $w_g^{(0)}$
\For{每个客户端 $k$ \textbf{in parallel}}
    \State 基于本地数据计算周期感知尺度分数 $\mathbf{s}_k$
    \State 按式~\eqref{eq:spilter-topm} 确定激活尺度集合 $\mathcal{R}_k$
\EndFor
\For{$t=0,1,\dots,T-1$}
    \For{每个客户端 $k$ \textbf{in parallel}}
        \State 从服务器拉取尺度参数 $\Theta_k^\downarrow(t)=\{\theta_g^{(r,t)}:r\in\mathcal{R}_k\}$
        \State 构造两种增强视图，按式~\eqref{eq:stitch} 将 $|\mathcal{R}_k|$ 个尺度特征拼接为子模型表示
        \State 在拼接表示上计算 local CL，在 $\mathcal{R}_k$ 各分量上计算 joint CL/KD 与 scale CL/KD
        \State 更新本地尺度参数并上传 $\Theta_k^\uparrow(t)$
    \EndFor
    \State 服务器按式~\eqref{eq:scale-agg} 对每个尺度分别聚合
    \State 得到全局模型 $w_g^{(t+1)}$
\EndFor
\end{algorithmic}
\end{algorithm}

\section{理论分析}

本节从表示信息、联邦漂移和系统效率三个角度解释 FedCSL-Spilter 的设计。本文不试图证明当前启发式尺度选择方案达到闭式全局最优，而是在温和假设下说明：全尺度 shapelet 训练存在信息冗余和系统拖尾；客户端\textbf{纯本地分布感知 top-$m$} 选择是在固定尺度预算下对有效多尺度信息的近似保留。

\subsection{多尺度 shapelet 表示的信息分解}

设原始时间序列为随机变量 $X$，两种随机增强视图为 $X^{(1)}$ 和 $X^{(2)}$。二者共享潜在语义因子 $C$，例如类别相关局部形态、主导周期、跨通道耦合或长期依赖；增强过程主要引入不改变语义的扰动因子。多尺度 shapelet 编码器输出
\begin{equation}
Z=\left[Z^{(1)},Z^{(2)},\ldots,Z^{(R)}\right],
\end{equation}
其中 $Z^{(r)}$ 表示第 $r$ 个尺度上的 shapelet 响应。InfoNCE 型局部对比目标可以视为最大化跨增强共享信息的可优化下界，因此多尺度局部目标可写为
\begin{equation}
\max_{\theta}\sum_{r=1}^{R}\alpha_r\,
I\!\left(Z^{(r)}(X^{(1)});Z^{(r)}(X^{(2)})\right).
\label{eq:mi-multiscale}
\end{equation}
式~\eqref{eq:mi-multiscale} 表明，多尺度 shapelet 并非简单堆叠特征，而是在不同时间跨度上分别保留增强不变语义：短尺度偏向局部瞬时模式，中尺度刻画过渡动态，长尺度保留周期与长期依赖。FedCSL-Spilter 的尺度选择问题也因此可以理解为：在不能训练全部 $R$ 个尺度时，如何保留对 $C$ 最有价值的尺度子集。

\subsection{尺度冗余与 Top-$m$ 信息保留}

全尺度训练的直觉优势是覆盖面广，但相邻尺度往往提取高度相似的局部片段。若第 $i$ 个尺度已经给出表征 $Z^{(i)}$，新增相邻尺度 $j$ 的边际判别信息可以写为条件互信息
\begin{equation}
\Delta_{j\mid i}=I\!\left(Z^{(j)};Y\mid Z^{(i)}\right).
\end{equation}
由于时间序列在相近窗口长度上通常具有连续性，邻近尺度之间的响应相关性较高，因而常出现 $\Delta_{j\mid i}\approx 0$ 的边际收益递减。训练全尺度模型会提升 $I(X;Z)$ 和本地模型容量，却不一定显著提升 $I(Z;Y)$；在每个客户端样本有限且分布非 IID 的联邦场景中，过大的本地容量还可能放大局部过拟合和客户端漂移。由此，Spilter 的 top-$m$ 尺度子集拼接可视为一种结构化正则化：它限制单个客户端每轮可拟合的尺度空间，同时通过不同客户端依据本地评分选择不同尺度组合形成多视图集成。

\subsection{周期评分与受限预算下的信息选择}

第~\ref{sec:period-score} 节中的 STFT--ACF 周期评分可以被看作尺度相关信息密度的代理。令 $P_k$ 表示客户端 $k$ 的周期/频谱结构因子，则尺度权重可解释为
\begin{equation}
\pi_{k,r}\propto I(P_k;Z_k^{(r)}).
\end{equation}
因此，加权多尺度损失
\begin{equation}
\sum_{r=1}^{R}\pi_{k,r}\mathcal{L}_{k,r}
\end{equation}
相当于把优化资源更多分配给与当前客户端主导周期更相关的尺度。

在 Spilter 中，客户端 $k$ 自主选择并只训练尺度子集 $\mathcal{R}_k=\mathrm{TopM}(\mathbf{s}_k,m)$。若尺度 $r$ 的有效收益为 $u_{k,r}\propto s_{k,r}$，则 top-$m$ 选择等价于在固定基数约束下最大化局部效用：
\begin{equation}
\max_{\mathcal{R}_k\subseteq \{1,\ldots,R\},|\mathcal{R}_k|=m}
\sum_{r\in\mathcal{R}_k}u_{k,r}.
\label{eq:budget-scale}
\end{equation}
若集合效用
\begin{equation}
F_k(\mathcal{R})=I(C_k;Z_{\mathcal{R}})
\end{equation}
近似满足单调性和边际收益递减，则按 $s_{k,r}$ 降序选取 top-$m$ 可视为子模最大化的贪心近似。该假设与多尺度 shapelet 的冗余现象相符：少量与本地周期高度匹配的尺度通常带来主要收益，继续加入低分相邻尺度的边际收益逐渐下降。由于每个客户端基于本地数据独立完成该选择，不同客户端的最终尺度组合天然不同，在全局层面形成多视图集成效应。

\subsection{非 IID 表示漂移与局部--全局对齐}

记客户端 $k$ 的局部目标为 $F_k(w)$，全局目标为
\begin{equation}
F(w)=\sum_{k=1}^{K}p_k F_k(w),\qquad p_k=\frac{n_k}{N}.
\end{equation}
若各客户端数据分布差异较大，则 $\nabla F_k(w)$ 与 $\nabla F(w)$ 之间的差异会显著增大。设存在异构上界 $\delta_k$ 使得
\begin{equation}
\|\nabla F_k(w)-\nabla F(w)\|\leq \delta_k,
\end{equation}
局部训练 $E$ 步后，客户端相对全局模型的漂移近似满足
\begin{equation}
\|w_k^{(t+1)}-w_g^{(t)}\|
\leq \eta E\left(\|\nabla F(w_g^{(t)})\|+\delta_k\right)
 + O(\eta^2E^2).
\end{equation}
这说明当 $\delta_k$ 较大时，仅依赖局部训练和参数平均会导致显著客户端漂移\cite{fedprox,moon}。对表示学习任务而言，这种漂移还会表现为各客户端嵌入空间的坐标系不一致\cite{fedu2,fedca}。

FedCSL-Spilter 中的 joint CL/KD 与 scale CL/KD 可解释为对表示分布散度的约束。设客户端表示分布为 $P_k(Z)$，全局教师表示分布为 $P_g(Z)$，蒸馏项近似减小 $D_{\mathrm{KL}}(P_k(Z)\|P_g(Z))$。若下游损失 $\ell(Z)$ 关于表示为 $L$-Lipschitz，则
\begin{equation}
\left|
\mathbb{E}_{P_k}\ell(Z)-\mathbb{E}_{P_g}\ell(Z)
\right|
\leq
L\cdot \mathrm{TV}\!\left(P_k(Z),P_g(Z)\right)
\leq
L\sqrt{\frac{1}{2}D_{\mathrm{KL}}\!\left(P_k(Z)\|P_g(Z)\right)} .
\end{equation}
因此，局部--全局蒸馏通过降低表示分布散度收缩客户端风险与全局风险之间的偏差上界。逐尺度蒸馏进一步避免某些尺度偏差被整体表示平均掩盖；结构约束则惩罚 shapelet 通道协方差非对角项，降低尺度内冗余，提升表示空间的可用性。

\subsection{系统效率、min--max 时延与误差界}

全尺度联邦 shapelet 学习的系统瓶颈来自同步训练中的最慢客户端。设客户端 $k$ 的计算速度为 $p_k$，其自主选择尺度集合的计算负载为
\begin{equation}
C_k=\sum_{r\in\mathcal{R}_k}F_r,\qquad
T_k^{\mathrm{comp}}=\frac{C_k}{p_k}.
\end{equation}
一轮同步训练的实际等待时间近似为
\begin{equation}
T_{\mathrm{round}}=
\max_{k\in\{1,\ldots,K\}}
\left(T_k^{\mathrm{comp}}+T_k^{\mathrm{comm}}\right).
\label{eq:round-makespan}
\end{equation}
若所有客户端都训练完整 $R$ 个尺度，则慢设备的 $C_k$ 与强设备相同；当 $p_k$ 很小时，式~\eqref{eq:round-makespan} 会由该客户端主导。Spilter 通过令每个客户端仅激活 $|\mathcal{R}_k|=m_k<R$ 个尺度来降低 $C_k$ 与通信量；在资源异构场景中，弱设备可进一步取更小的 $m_k$（见 §\ref{sec:resource-hetero}），而\textbf{具体激活哪些尺度仍由本地周期评分 top-$m_k$ 决定}，不在评分函数中写入系统成本代理。

另一方面，Spilter 可以视为对全尺度目标的结构化子采样。设
\begin{equation}
\mathcal{L}_{k}^{\mathrm{full}}=\sum_{r=1}^{R}\alpha_{k,r}\mathcal{L}_{k,r},
\qquad
\mathcal{L}_{k}^{\mathrm{split}}=\sum_{r\in\mathcal{R}_k}\tilde{\alpha}_{k,r}\mathcal{L}_{k,r}.
\end{equation}
若每个尺度损失有界 $0\leq\mathcal{L}_{k,r}\leq L_{\max}$，则有
\begin{equation}
\left|\mathcal{L}_{k}^{\mathrm{full}}-\mathcal{L}_{k}^{\mathrm{split}}\right|
\leq
L_{\max}\sum_{r\notin\mathcal{R}_k}\alpha_{k,r}
 + \Delta_{\mathrm{renorm}},
\label{eq:split-error-bound}
\end{equation}
其中 $\Delta_{\mathrm{renorm}}$ 为已选尺度重新归一化引入的误差。周期评分 top-$m$ 的作用正是优先保留高权重尺度，减小式~\eqref{eq:split-error-bound} 右侧第一项。

从通信角度看，定义通信加速倍率
\begin{equation}
G_{\mathcal{C}}(t)=
\frac{\mathcal{C}_{\mathrm{full}}(t)}
{\mathcal{C}_{\mathrm{split}}(t)}
\approx
\frac{K\sum_{r=1}^{R}B_r}
{\sum_{k=1}^{K}\sum_{r\in\mathcal{R}_k}B_r},
\end{equation}
则通信节省比例为 $1-1/G_{\mathcal{C}}(t)$。若各尺度通信成本近似均匀，且每个客户端激活 $m$ 个尺度，则
\begin{equation}
G_{\mathcal{C}}(t)\approx \frac{R}{m},
\qquad
1-\frac{1}{G_{\mathcal{C}}(t)}\approx 1-\frac{m}{R}.
\end{equation}
计算加速同理为
\begin{equation}
G_{\mathcal{F}}(t)=
\frac{\mathcal{F}_{\mathrm{full}}(t)}
{\mathcal{F}_{\mathrm{split}}(t)}
=
\frac{\sum_{k=1}^{K}n_k\sum_{r=1}^{R}F_r}
{\sum_{k=1}^{K}n_k\sum_{r\in\mathcal{R}_k}F_r}.
\end{equation}
当长尺度计算更重时，随机 top-$m$ 与本地评分 top-$m$ 的平均 $\sum_{r\in\mathcal{R}_k}F_r$ 可能不同；本文不在选择阶段显式优化该成本，而将其作为资源异构实验中的观测指标（§\ref{sec:resource-hetero}）。综上，FedCSL-Spilter 的统一目标可概括为
\begin{equation}
\max_{\{\mathcal{R}_k\},\theta}
\sum_{k=1}^{K}
\sum_{r\in\mathcal{R}_k}
\pi_{k,r}
I\!\left(Z_k^{(r,1)};Z_k^{(r,2)}\right)
-\lambda D\!\left(P_k(Z)\|P_g(Z)\right)
-\lambda_c\mathcal{C}_{\mathrm{split}}
-\lambda_f\mathcal{F}_{\mathrm{split}},
\end{equation}
其中 $\mathcal{B}_k\subseteq\mathcal{R}_k$，$|\mathcal{R}_k|=m$，且每个 $\mathcal{R}_k$ 由客户端 $k$ 基于本地数据自主确定。该形式体现了本文方法的核心折中：在非 IID 联邦环境中，各客户端独立识别并保留对自身最有价值的多尺度语义信息，同时通过局部--全局对齐和覆盖平衡惩罚维持全局表示空间的一致性。

\section{实验设计}

\subsection{数据集与联邦划分}

本文全部实验\textbf{仅}在以下五个多变量时间序列分类数据集上进行，不再引入其他 UEA 基准或额外扩展数据集：

\begin{itemize}[leftmargin=2em]
    \item \textbf{HAR}（UCI Human Activity Recognition）：智能手机惯性传感器采集的 6 类人体活动识别数据，序列长度 128、9 个传感通道，样本规模较大且动作周期结构清晰。用于表~\ref{tab:har_dirichlet} 的多 $\alpha$ 非 IID 对比、图~\ref{fig:dirichlet} 收敛曲线，以及表~\ref{tab:har_client_memory}--\ref{tab:har_local_compute} 的系统效率实测；并在表~\ref{tab:performance_comparison_rounded} 中报告消融结果。
    \item \textbf{FaceDetection}（UEA）：多变量人脸检测序列，长度 62、3 通道、2 类，序列较短、类别均衡性较弱。用于表~\ref{tab:scale_performance} 的多 $\alpha$ 异构对比。
    \item \textbf{LSST}（UEA）：天文光变曲线分类数据，长度 36、6 通道、14 类，具有典型的周期性天文观测模式。用于表~\ref{tab:lsst_dirichlet} 的多 $\alpha$ 对比。
    \item \textbf{Epilepsy}（UEA）：癫痫相关 EEG 风格序列，长度 206、3 通道、4 类。用于表~\ref{tab:performance_comparison_rounded} 的模块消融。
    \item \textbf{SleepEDF}：Sleep-EDF 睡眠分期数据，由多通道 EEG/EOG 构成多变量序列、5 类睡眠阶段（W/N1/N2/N3/REM），序列较长且跨客户端分布差异显著。用于表~\ref{tab:performance_comparison_rounded} 的模块消融实验；资源异构场景分析（表~\ref{tab:resource_hetero}）亦以 HAR 为主数据集展开。
\end{itemize}

为模拟真实联邦环境，FaceDetection、LSST 与 HAR 的主实验在多个 Dirichlet 参数 $\alpha$ 下重复；Epilepsy 与 SleepEDF 的消融实验采用默认非 IID 划分。设客户端数量为 $K$，Dirichlet 参数为 $\alpha$，则较小的 $\alpha$ 对应更强的数据异构性。

\subsection{对比方法}

本文从两个层面设置对比方法。第一类是\textbf{联邦通用学习方法与现有时序表示学习基线}。联邦侧选取 FedAvg\cite{fedavg}、FedProx\cite{fedprox}、MOON\cite{moon}、FedX\cite{fedx}、FedProto\cite{fedproto}、SCAFFOLD\cite{scaffold} 等代表性方法，用于比较不同全局聚合、近端约束、模型级对比、原型对齐和控制变量校正机制在非 IID 时序数据上的表现。时序表示侧选取 TS2Vec\cite{ts2vec}、TS-TCC\cite{tstcc}、TF-C\cite{tfc}、TimesNet\cite{timesnet} 以及 TimeCSL/CSL\cite{timescsl} 等方法，覆盖 Transformer/CNN 主干、时间--频率对比学习和 shapelet-based contrastive representation 等不同路线。对于集中式时序表示方法，实验中可将其作为非联邦表示上界或在各客户端本地训练后进行统一评估，从而观察联邦约束带来的性能变化。

第二类是\textbf{FedAvg + 本文 shapelet 骨干}的结构化基线。该组方法保持本文的多尺度 shapelet 编码器和相同下游 SVM 评估协议，但仅采用 FedAvg 式参数平均，不引入局部--全局 joint CL/KD、尺度级 CL/KD、周期感知加权和 Spilter 本地 top-$m$ 拼接。该设置用于隔离”shapelet 表示本身”与”本文联邦对齐和尺度拼接机制”的贡献。除主实验外，本文保留 Full-scale FedCSL、Random-top-$m$ Spilter 与 Local-score-top-$m$ Spilter 等消融版本，用于分析本地分布感知尺度选择与随机 top-$m$ 的差异。

\subsection{训练设置}

每轮训练开始前，各客户端首先根据本地数据的周期/频谱特征按式~\eqref{eq:spilter-topm} 确定激活尺度集合，从服务器拉取对应尺度参数后进行本地训练。客户端本地训练若干 epoch 后只上传被训练尺度，服务器按式~\eqref{eq:scale-agg} 完成尺度级聚合。模型主干采用多尺度 shapelet 编码器，距离度量可设为混合形式；优化器采用 SGD；温度参数、损失系数、客户端数量、通信轮数、本地 epoch 数、批大小等均由配置文件统一控制。若启用 Spilter，应在超参数表中单独列出每客户端激活尺度数 $m$（或资源异构场景下的 $\{m_k\}$）及 \texttt{allocation\_mode=local\_score\_topm}。

\subsection{评估协议}

本文不直接使用全局模型进行端到端分类，而是先利用全局 shapelet 编码器提取训练集和测试集表示，再在表示空间上训练 SVM 分类器。此设置与 TimeCSL 的“representation + task-oriented analyzer”范式一致\cite{timescsl}，也与 TS2Vec/TSTCC 等自监督时序表示工作的线性评测协议相容\cite{ts2vec,tstcc}，更能反映所学习表示本身的质量。各方法统一报告训练\textbf{最后 10 轮}测试集 SVM 平均准确率；评估指标还包括平均排名、收敛轮数以及可选的训练时间指标；若后续扩展到聚类或异常检测任务，还可引入 NMI、RI 或 F1-score 等指标。

\section{结果分析与讨论}

\subsection{总体性能对比}

总体实验应重点比较以下方法：集中式 TimeCSL/CSL、FedAvg/FedCSL 全尺度训练、Random-top-$m$ Spilter 与 Local-score-top-$m$ Spilter。预期结果是，集中式方法在数据可集中时仍然具有较强上界表现；全尺度联邦训练通常具有较强表示能力但通信与计算开销较高；Local-score-top-$m$ Spilter 在只训练 $m$ 个尺度的条件下，通过本地周期评分选取最有价值的尺度组合，力求接近全尺度性能，同时显著降低每轮通信量与本地计算量。表~\ref{tab:scale_performance}、表~\ref{tab:lsst_dirichlet} 与表~\ref{tab:har_dirichlet} 分别汇总 FaceDetection、LSST 与 HAR 在不同 Dirichlet 非 IID 强度下的 SVM 测试精度；表~\ref{tab:performance_comparison_rounded} 则在 Epilepsy、HAR 与 SleepEDF 上报告模块消融结果。

\begin{table*}[htbp]
  \centering
  \caption{FaceDetection 数据集在不同 Dirichlet 非 IID 强度下的算法准确率对比（评估指标：最后 10 轮测试集 SVM 平均准确率；加粗为该列最优）。}
  \label{tab:scale_performance}
  \resizebox{\textwidth}{!}{%
  \begin{tabular}{lcccccc}
    \toprule
    Method & $\alpha=0.05$ & $\alpha=0.1$ & $\alpha=0.3$ & $\alpha=0.5$ & $\alpha=1.0$ & $\alpha=10.0$ \\
    \midrule
    FedAvg            & 0.5594          & 0.5633          & 0.5459          & 0.5386          & 0.5501          & 0.5376          \\
    FedProx           & 0.5430          & 0.5438          & 0.5486          & 0.5508          &                 & 0.5336          \\
    FedBYOL           & 0.5530          & 0.5592          & 0.5627          & 0.5616          & 0.5763          & 0.5578          \\
    FedU2             & 0.5621          & 0.5634          & 0.5463          & 0.5630          & 0.5550          & 0.5618          \\
    Orchestra         & 0.5609          & 0.5530          & 0.5646          & 0.5620          & 0.5591          & \textbf{0.5632} \\
    FedCSL            & \textbf{0.5710} & \textbf{0.5646} & \textbf{0.5701} & \textbf{0.5737} & \textbf{0.5770} & 0.5626          \\
    \midrule
    Spilter ($m{=}1$) &                 &                 &                 &                 &                 &                 \\
    Spilter ($m{=}2$) &                 &                 &                 &                 &                 &                 \\
    Spilter ($m{=}4$) &                 &                 &                 &                 &                 &                 \\
    \bottomrule
  \end{tabular}%
  }
\end{table*}

\begin{table*}[htbp]
  \centering
  \caption{LSST 数据集在不同 Dirichlet 非 IID 强度下的算法准确率对比（评估指标：最后 10 轮测试集 SVM 平均准确率；加粗为该列最优）。}
  \label{tab:lsst_dirichlet}
  \resizebox{\textwidth}{!}{%
  \begin{tabular}{lcccccc}
    \toprule
    Method & $\alpha=0.05$ & $\alpha=0.1$ & $\alpha=0.3$ & $\alpha=0.5$ & $\alpha=1.0$ & $\alpha=10.0$ \\
    \midrule
    FedAvg            & 0.5857          & 0.5833          & 0.5901          & 0.5716          & 0.5770          & 0.5739          \\
    FedProx           & 0.5867          & 0.5951          & 0.5793          & 0.5777          & 0.5781          & 0.5794          \\
    FedBYOL           & 0.6033          & 0.6092          & 0.6071          & 0.6056          & 0.6051          & 0.6025          \\
    FedU2             & \textbf{0.6103} & \textbf{0.6102} & \textbf{0.6115} & \textbf{0.6185} & \textbf{0.6201} & \textbf{0.6138} \\
    Orchestra         & 0.6068          & 0.6066          & 0.6084          & 0.6095          & 0.6020          & 0.6018          \\
    FedCSL            & 0.5978          & 0.6054          & 0.5992          & 0.5926          & 0.6008          & 0.5968          \\
    \midrule
    Spilter ($m{=}4$) & 0.6013          & 0.5995          & 0.5897          & 0.6052          & 0.5940          & 0.6081          \\
    \bottomrule
  \end{tabular}%
  }
\end{table*}

\begin{table}[htbp]
    \centering
    \caption{不同方法在各个数据集上的性能对比 (消融实验)}
    \label{tab:performance_comparison_rounded}
    \begin{tabular}{lcccc}
        \toprule
        \multirow{2}{*}{Dataset} & \multirow{2}{*}{local } & \multirow{2}{*}{ AVG} & \multirow{2}{*}{ 联邦对比All} &  联邦对比 \\
        & & & & +多尺度自适应 \\
        \midrule
        Epilepsy & & 0.9804 & 0.9839 & 0.9852 \\
        HAR      & & 0.9281 & 0.9342 & 0.9338 \\
        SleepEDF & & 0.7797 & 0.8052 & 0.8213 \\
        \bottomrule
    \end{tabular}
\end{table}

\begin{table*}[htbp]
  \centering
  \caption{HAR 数据集在不同 Dirichlet 非 IID 强度下的算法准确率对比（评估指标：最后 10 轮测试集 SVM 平均准确率）。加粗为该列最优。}
  \label{tab:har_dirichlet}
  \resizebox{\textwidth}{!}{%
  \begin{tabular}{lcccccc}
    \toprule
    Method & $\alpha=0.05$ & $\alpha=0.1$ & $\alpha=0.3$ & $\alpha=0.5$ & $\alpha=1.0$ & $\alpha=10.0$ \\
    \midrule
    FedAvg            & 0.9332          & 0.9370          & 0.9338          & 0.9462          & 0.9445          & 0.9484          \\
    FedProx           & \textbf{0.9472} & 0.9399          & 0.9346          & 0.9384          & 0.9398          & 0.9458          \\
    FedBYOL           & 0.9268          & 0.9285          & 0.9336          & 0.9284          & 0.9331          & 0.9267          \\
    FedU2             & 0.9252          & 0.9240          & 0.9307          & 0.9286          & 0.9317          & 0.9327          \\
    Orchestra         & 0.9344          & 0.9336          & 0.9293          & 0.9315          & 0.9325          & 0.9333          \\
    FedCSL            & 0.9315          & 0.9339          & 0.9306          & 0.9382          & 0.9475          & 0.9514          \\
    \midrule
    Spilter ($m{=}1$) & 0.9261          & 0.9236          & 0.9350          & 0.9378          & 0.9253          & 0.9331          \\
    Spilter ($m{=}2$) & 0.9377          & 0.9346          & 0.9489          & 0.9527          & 0.9405          & \textbf{0.9655} \\
    Spilter ($m{=}4$) & 0.9414          & \textbf{0.9446} & \textbf{0.9521} & \textbf{0.9586} & \textbf{0.9573} & 0.9630          \\
    \bottomrule
  \end{tabular}%
  }
\end{table*}
\subsection{消融实验}

消融实验用于回答”性能提升来自哪里”这一关键问题。具体而言，可依次关闭以下模块：joint CL、joint KD、scale CL、scale KD、ACF 周期感知加权和本地分布感知 top-$m$ 选择（改为 Random-top-$m$）。若关闭联合蒸馏后精度明显下降，则说明在非 IID 条件下特征层面对齐确实能够有效抑制客户端漂移；若关闭周期感知加权后在 Epilepsy、HAR 与 SleepEDF 上的平均精度下降，则说明不同尺度的区分性并不相同，统一等权并不是最优选择；若将本地评分 top-$m$ 替换为随机 top-$m$，则验证周期感知评分对尺度选择的收益。

\subsection{尺度冗余性与多尺度互补性}

为验证式~\eqref{eq:mi-multiscale} 所述的尺度冗余和多尺度互补性，可在集中式或服务器侧评估环境下分别训练单尺度 shapelet 编码器，提取测试集表示后计算尺度间余弦相似度、互信息估计或 CKA 相似度。若相邻尺度相似度显著高于短尺度与长周期尺度之间的相似度，则说明全尺度模型存在边际信息冗余；客户端按本地周期评分 top-$m$ 选取的尺度组合，则倾向于保留与本地主导周期最匹配的互补频段。

\subsection{边际效用与结构化正则分析}

为进一步验证”少量互补尺度即可逼近全尺度性能”的结论，可固定联邦划分和训练预算，逐步增加每客户端激活尺度数 $m$，并记录下游精度、参数量、FLOPs、通信量和相对全尺度精度。若精度在 $m=4$ 或 $m=5$ 附近进入平台期，而计算和通信继续显著上升，则支持尺度收益的边际递减假设；若客户端按本地评分拼接的尺度子模型在强 non-IID 下超过全尺度训练，则可进一步说明 Spilter 的自适应尺度拼接具有结构化正则化效果（参见表~\ref{tab:har_dirichlet} 与表~\ref{tab:har_client_memory}）。

\subsection{非 IID 强度与收敛趋势分析}
\begin{figure}[htbp]
    \centering
    \includegraphics[width=0.95\linewidth]{../../figs/har_dirichlet_convergence.png}
    \caption{HAR 数据集上不同 Dirichlet 参数 $\alpha$ 下各联邦学习方法的测试准确率随通信轮次变化曲线（$3{\times}2$ 子图：自上而下、自左而右对应 $\alpha=0.05,0.1,0.3,0.5,1.0,10.0$；Spilter 按 \texttt{local\_top\_m}$\in\{1,2,4\}$ 画为三条曲线；各子图纵轴按该 $\alpha$ 下曲线范围自适应放大）。数据来自 \texttt{data/HAR/alpha\_*} 真实训练日志。}
    \label{fig:dirichlet}
\end{figure}
通过改变 Dirichlet 参数 $\alpha$（图~\ref{fig:dirichlet} 给出各 $\alpha$ 下 round 维度的收敛曲线；表~\ref{tab:har_dirichlet} 汇总最后 10 轮平均精度），可以分析各方法在不同异构程度下的表现。当 $\alpha$ 较小时，客户端数据分布差异更大，此时若 Spilter 和局部--全局蒸馏机制依然能保持较好性能，则说明本文方法对联邦异构性具有较强鲁棒性。

\subsection{异构系统下的真实时间效率}

除以通信轮数为横轴的收敛曲线外，还应在异构设备或模拟异构速度下报告真实挂钟时间与\textbf{客户端显存占用}。表~\ref{tab:har_client_memory} 基于 \texttt{data/system\_efficiency\_HAR\_partials} 中的 HAR 实测结果：在 Dirichlet $\alpha=0.1$、$K=10$、batch size $=32$、本地 1 epoch 的设置下，对每个客户端串行独占 GPU 训练，在 warmup 后调用 \texttt{torch.cuda.reset\_peak\_memory\_stats()}，以 \texttt{max\_memory\_allocated()} 统计峰值显存，再对全部客户端取均值与最大值。FedCSL 等价于激活全部 $R=8$ 个尺度；Spilter-$m$ 仅对所选 $m$ 个尺度做 stitched 子模型前向，因此激活显存随 $m$ 近似线性下降。

\begin{table*}[htbp]
  \centering
  \caption{HAR 客户端本地训练峰值显存与最慢客户端单 epoch 耗时（Dirichlet $\alpha=0.1$，RTX 3090，$K=10$）。Mean / Max Peak Mem 为各客户端 1-epoch 峰值显存的均值 / 最大值；Mem Saving 与 Compress.\ Ratio 均以 FedCSL 为基准：Saving $=1-\bar M/\bar M_{\mathrm{FedCSL}}$，Compress.\ $=\bar M_{\mathrm{FedCSL}}/\bar M$。加粗为 Spilter 相对 FedCSL 的显存节省。}
  \label{tab:har_client_memory}
  \resizebox{\textwidth}{!}{%
  \begin{tabular}{lcccccc}
    \toprule
    Method & $m$ & Mean Peak Mem (MB) $\downarrow$ & Max Peak Mem (MB) $\downarrow$ & Mem Saving $\uparrow$ & Compress.\ Ratio $\uparrow$ & Slowest Epoch (s) $\downarrow$ \\
    \midrule
    FedAvg            & 8 & 177.0 & 180.3 & 0.6\%  & 1.01$\times$ & \textbf{1.28} \\
    FedProx           & 8 & 179.2 & 182.4 & $-$0.6\% & 0.99$\times$ & 1.47 \\
    FedBYOL           & 8 & 182.9 & 183.4 & $-$2.7\% & 0.97$\times$ & 1.89 \\
    FedU2             & 8 & 182.9 & 183.4 & $-$2.7\% & 0.97$\times$ & 1.93 \\
    Orchestra         & 8 & 231.3 & 233.5 & $-$29.9\% & 0.77$\times$ & 2.28 \\
    FedCSL (baseline) & 8 & 178.1 & 181.3 & ---    & 1.00$\times$ (ref) & 4.70 \\
    \midrule
    \textbf{Spilter-$m{=}1$} & 1 & \textbf{51.5} & \textbf{51.7} & \textbf{71.1\%} & \textbf{3.46$\times$} & 3.17 \\
    \textbf{Spilter-$m{=}2$} & 2 & \textbf{78.7} & \textbf{79.1} & \textbf{55.8\%} & \textbf{2.26$\times$} & 3.43 \\
    \textbf{Spilter-$m{=}4$} & 4 & \textbf{130.3} & \textbf{131.5} & \textbf{26.8\%} & \textbf{1.37$\times$} & 3.76 \\
    \bottomrule
  \end{tabular}%
  }
\end{table*}

表~\ref{tab:har_client_memory} 表明：在保持与 FedCSL 相同 shapelet 骨干与多尺度 CL/KD 训练路径的前提下，\textbf{仅通过尺度子模型拼接}即可显著降低客户端峰值显存——Spilter-$m{=}1$ 将平均峰值从 178.1\,MB 降至 51.5\,MB（节省 71.1\%，约 3.46$\times$ 压缩）；Spilter-$m{=}4$ 仍节省 26.8\% 显存，且在表~\ref{tab:har_dirichlet} 的多数 $\alpha$ 下精度优于或接近 FedCSL，形成精度--显存的 Pareto 优势。相比之下，FedAvg / FedProx 虽单 epoch 更快，但显存与 FedCSL 同级；Orchestra 因额外表征分支峰值显存最高（233.5\,MB）。Spilter 的最慢客户端 epoch 仍低于 FedCSL（3.17--3.76\,s vs.\ 4.70\,s），与式~\eqref{eq:round-makespan} 中``降低 $C_k$ 以缩短 min--max 时延''的分析一致。

表~\ref{tab:har_local_compute} 进一步给出 HAR 上各方法的\textbf{本地训练算力下界}：在 Dirichlet $\alpha=0.1$、$K=10$、batch size $=32$、$E=1$、$T=200$ 下，用 \texttt{system\_efficiency\_HAR\_partials} 测得的 $\max_k T_k^{\mathrm{epoch}}$ 估计同步联邦每轮本地阶段的最慢客户端耗时，并以 $T_k^{\mathrm{epoch,max}}\times T$ 给出 200 轮\textbf{纯本地计算}累计下界（不含通信、聚合与 SVM 评测）。该 $\alpha=0.1$ 划分即 Spilter 取得最大 Epoch Speedup 的实测设置（Spilter-$m{=}1$ 达 1.48$\times$，为 $m\in\{1,2,4\}$ 中最高）；表中 Final Acc.\ 与 Epoch Speedup 均取该 $\alpha$ 下表~\ref{tab:har_dirichlet} 的最后 10 轮精度与微基准耗时。

\begin{table*}[htbp]
  \centering
  \caption{HAR 客户端本地训练耗时与 200 轮计算下界（Dirichlet $\alpha=0.1$，RTX 3090，$K=10$）。Slowest / Mean Epoch 与 Epoch Speedup 来自 \texttt{system\_efficiency\_HAR\_partials}；200-round Local LB $=\max_k T_k^{\mathrm{epoch}}\times 200$（分钟）；Final Acc.\ 为同 $\alpha$ 下表~\ref{tab:har_dirichlet} 的最后 10 轮精度。该 $\alpha=0.1$ 为当前系统效率实测划分，Spilter-$m{=}1$ 的 Epoch Speedup（1.48$\times$）为 Spilter 三者最高。}
  \label{tab:har_local_compute}
  \resizebox{\textwidth}{!}{%
  \begin{tabular}{lcccccc}
    \toprule
    Method & Slowest Epoch (s) $\downarrow$ & Mean Epoch (s) & 200-round Local LB (min) $\downarrow$ & Final Acc.\ ($\alpha{=}0.1$) $\uparrow$ & Epoch Speedup $\uparrow$ \\
    \midrule
    FedAvg            & \textbf{1.28} & 0.52 & 4.3  & 0.9370 & 3.66$\times$ \\
    FedProx           & 1.47 & 0.64 & 4.9  & 0.9399 & 3.20$\times$ \\
    FedBYOL           & 1.89 & 0.88 & 6.3  & 0.9285 & 2.49$\times$ \\
    FedU2             & 1.93 & 0.92 & 6.4  & 0.9240 & 2.43$\times$ \\
    Orchestra         & 2.28 & 1.10 & 7.6  & 0.9336 & 2.06$\times$ \\
    FedCSL (baseline) & 4.70 & 2.05 & 15.7 & 0.9339 & 1.00$\times$ \\
    \midrule
    Spilter-$m{=}1$   & 3.17 & 1.34 & 10.6 & 0.9236 & \textbf{1.48$\times$} \\
    Spilter-$m{=}2$   & 3.43 & 1.44 & 11.4 & 0.9346 & 1.37$\times$ \\
    Spilter-$m{=}4$   & 3.76 & 1.62 & 12.5 & \textbf{0.9446} & 1.25$\times$ \\
    \bottomrule
  \end{tabular}%
  }
\end{table*}

FedAvg 系列本地 epoch 最短，但在 Spilter Epoch Speedup 最优的 $\alpha=0.1$ 划分下，Spilter-$m{=}4$ 精度（0.945）仍高于 FedCSL（0.934）；Spilter-$m{=}1$ 加速最高（1.48$\times$），200 轮本地下界 10.6 分钟，约为 FedCSL（15.7 分钟）的 68\%。

\subsection{系统资源不平衡场景下的效果分析}
\label{sec:resource-hetero}

§\ref{sec:spilter} 中的 Spilter 尺度\textbf{排序}仅依赖本地 STFT/ACF 周期评分；但在真实联邦部署里，$K=10$ 个客户端往往同时存在显存与算力差异：弱设备无法承载全尺度 FedCSL，甚至无法使用与大设备相同的 $m$。本节因此考察\textbf{设备侧资源预算}与\textbf{数据侧本地 top-$m$ 选择}解耦后的系统行为：先由硬件档位决定每客户端可用的尺度预算 $m_k$，再在该预算内按式~\eqref{eq:spilter-topm} 选取具体尺度并拼接子模型。

\paragraph{异构客户端设定。}
在 HAR、Dirichlet $\alpha=0.1$、同步联邦与 $K=10$ 的设置下，假定客户端按峰值显存分为三档（数值待实测填充）：
\begin{itemize}[leftmargin=2em]
    \item \textbf{小显存设备}（4 台）：峰值显存预算 $\leq 64$\,MB，仅允许 $m_k=1$；
    \item \textbf{中显存设备}（3 台）：峰值显存预算 $\approx 128$\,MB，允许 $m_k=2$；
    \item \textbf{大显存设备}（3 台）：峰值显存预算 $\geq 256$\,MB，允许 $m_k=4$ 或全尺度 $m_k=R$。
\end{itemize}
各档设备在本地评分确定 $\mathcal{R}_k$ 后，仍参与同一全局尺度级聚合；弱设备因 $m_k$ 较小而降低单轮通信量与本地 epoch 耗时，但可能遗漏部分尺度上的梯度贡献。对比基线包括：（1）\textbf{FedCSL 全尺度}——仅大显存设备可参与或弱设备因 OOM 退出；（2）\textbf{Uniform $m{=}2$}——所有设备统一 $m_k=2$；（3）\textbf{Spilter 资源自适应}——按上表配置 $\{m_k\}$，尺度集合仍由本地评分 top-$m_k$ 决定。

表~\ref{tab:resource_hetero} 汇总该场景下的参与率、单轮最慢客户端耗时、200 轮本地计算下界与全局 SVM 精度。当前数值留空，待在同一 \texttt{system\_efficiency} 脚本下按客户端档位注入显存上限后补测。

\begin{table*}[htbp]
  \centering
  \caption{系统资源不平衡场景下的联邦训练效果（HAR，Dirichlet $\alpha=0.1$，$K=10$；小/中/大显存设备分别为 4/3/3 台）。Participation 为可完成本地 epoch 的客户端比例；Slowest Epoch 为同步轮次最慢客户端耗时；Final Acc.\ 为全局表示 + SVM 测试精度。数值待补。}
  \label{tab:resource_hetero}
  \resizebox{\textwidth}{!}{%
  \begin{tabular}{lcccccc}
    \toprule
    Method & Small ($m_k{=}1$) & Medium ($m_k{=}2$) & Large ($m_k{=}4$) & Participation $\uparrow$ & Slowest Epoch (s) $\downarrow$ & Final Acc.\ $\uparrow$ \\
    \midrule
    FedCSL (full-scale only) & --- & --- & --- & --- & --- & --- \\
    Uniform Spilter ($m_k{=}2$) & --- & --- & --- & --- & --- & --- \\
    Spilter (adaptive $\{m_k\}$) & --- & --- & --- & --- & --- & --- \\
    \bottomrule
  \end{tabular}%
  }
\end{table*}

预期观察如下：（1）全尺度 FedCSL 在严格显存约束下参与率最低，同步轮次仍受最慢大设备主导；（2）资源自适应 Spilter 允许全部 10 台设备参与，并通过减小弱设备的 $m_k$ 降低式~\eqref{eq:round-makespan} 中的拖尾；（3）由于具体尺度仍由本地周期评分决定，资源自适应 Spilter 应比 Uniform 随机/固定 $m$ 更能保留与本地数据匹配的尺度组合，从而在相近系统开销下获得更高 Final Acc.。

\subsection{可解释性分析}

与纯 CNN 或 Transformer 表示相比，FedCSL 的一个重要优势在于 shapelet 表示的可解释性。实验中可以进一步展示：不同类别样本在表示空间中的降维分布，不同尺度 shapelet 的响应强度，以及周期感知加权下各尺度权重的变化趋势。通过这些分析，可以更直观地说明本文方法不仅提升了精度，也保留了模型的可解释性和可分析性。

\subsection{局限性讨论}

尽管本文提出了 Spilter 自适应尺度拼接与联邦 shapelet 表征学习的联合框架，但仍存在若干局限。第一，§\ref{sec:resource-hetero} 的资源异构实验尚未完成实测填充，当前仅给出 10 客户端三档显存设定与评价指标。第二，本文当前实验重点仍在分类任务，聚类与异常检测等下游任务还需进一步扩展\cite{ts2vec,timescsl}。第三，客户端尺度选择目前主要在训练前基于本地数据统计预计算并固定，未来可进一步设计随训练轮次动态更新的在线 top-$m$ 策略，使尺度组合更好地响应模型训练阶段和客户端数据分布变化。

\section{结论}

本文围绕联邦非 IID 多变量时间序列场景中的”表示学习”与”自适应尺度组合”两个核心问题，提出了 FedCSL-Spilter 框架。该方法以 TimeCSL/CSL 的多尺度 shapelet 表征学习为骨干，在联邦环境下进一步引入客户端自主尺度拼接机制：各客户端根据本地 STFT/ACF 周期评分按 top-$m$ 选取尺度并将其 shapelet 特征拼接为本地子模型进行训练，仅上传被训练尺度的参数，服务器按尺度粒度聚合。同时，本文引入局部--全局联合对比与知识蒸馏、多尺度结构对齐以及周期感知加权机制；设备资源差异则通过为不同客户端配置不同的尺度预算 $m_k$ 来适配（§\ref{sec:resource-hetero}），而不在尺度排序中混入系统成本项。理论分析表明，所提方法能够从特征层面抑制客户端漂移，并从系统层面降低尺度冗余带来的通信和计算开销；在 HAR、FaceDetection、LSST、Epilepsy 与 SleepEDF 上的实验则从总体性能、消融实验、非 IID 鲁棒性和可解释性等多个角度验证其有效性。未来工作将补全资源异构场景实测，并完善在线 top-$m$ 策略与个性化联邦扩展。


\begin{ack}
本节可在正式投稿前补充项目来源、基金支持、实验平台与协作者贡献等信息。匿名投稿时请勿写入可识别信息；终稿需按 NeurIPS 要求披露资助与利益冲突。
\end{ack}

% References (natbib is loaded by neurips_2023)
{\small
\bibliographystyle{unsrtnat}
\bibliography{refs}
}

\end{document}
